% Generated by `taskgen course examples` from number_theory.crt (seed 1).
% Do not edit: change examples.yaml and regenerate.
\begin{example}[система с попарно взаимно простыми модулями]\label{ex:crt-coprime}
Решить систему сравнений с помощью китайской теоремы об остатках (обратные элементы $y_i = M_i^{-1} \bmod m_i$ найти расширенным алгоритмом Евклида):\[\left\{\begin{array}{l}x \equiv 1\ (\textrm{mod}\ 3) \\ x \equiv 11\ (\textrm{mod}\ 14) \\ x \equiv 20\ (\textrm{mod}\ 23)\end{array}\right.\]

\par\noindent\textit{Решение.} Модули попарно взаимно просты: $\gcd(3, 14) = 1$, $\gcd(3, 23) = 1$, $\gcd(14, 23) = 1$; поэтому решение существует и единственно по модулю $M = 3 \cdot 14 \cdot 23 = 966$.\par\noindent $i = 1$: $M_{1} = M/3 = 322 \equiv 1\ (\textrm{mod}\ 3)$, $y_{1} = 1$ (очевидно: $1 \cdot 1 \equiv 1$).\par\noindent $i = 2$: $M_{2} = M/14 = 69 \equiv 13\ (\textrm{mod}\ 14)$.  Обратный к $13$ по модулю $14$ (расширенный алгоритм Евклида для $(14, 13)$):
\[\begin{array}{r|r|r|r}
r_i & q_i & u_i & v_i \\ \hline
14 &  & 1 & 0 \\
13 &  & 0 & 1 \\
1 & 1 & 1 & -1 \\
0 & 13 & -13 & 14 \\
\end{array}\qquad r_i = u_i \cdot 14 + v_i \cdot 13\]

$13^{-1} \equiv -1 \equiv 13\ (\textrm{mod}\ 14)$.
 Итак, $y_{2} = 13$ (проверка: $13 \cdot 13 = 169 \equiv 1$).\par\noindent $i = 3$: $M_{3} = M/23 = 42 \equiv 19\ (\textrm{mod}\ 23)$.  Обратный к $19$ по модулю $23$ (расширенный алгоритм Евклида для $(23, 19)$):
\[\begin{array}{r|r|r|r}
r_i & q_i & u_i & v_i \\ \hline
23 &  & 1 & 0 \\
19 &  & 0 & 1 \\
4 & 1 & 1 & -1 \\
3 & 4 & -4 & 5 \\
1 & 1 & 5 & -6 \\
0 & 3 & -19 & 23 \\
\end{array}\qquad r_i = u_i \cdot 23 + v_i \cdot 19\]

$19^{-1} \equiv -6 \equiv 17\ (\textrm{mod}\ 23)$.
 Итак, $y_{3} = 17$ (проверка: $19 \cdot 17 = 323 \equiv 1$).\par По формуле $x \equiv \sum_i r_i M_i y_i \pmod M$:\[x \equiv 1 \cdot 322 \cdot 1 + 11 \cdot 69 \cdot 13 + 20 \cdot 42 \cdot 17 = 24469 \equiv 319\ (\textrm{mod}\ 966).\]Проверка: $319 = 3 \cdot 106 + 1$; $319 = 14 \cdot 22 + 11$; $319 = 23 \cdot 13 + 20$.\[x \equiv 319\ (\textrm{mod}\ 966).\]
\end{example}
